\documentclass[a4paper,landscape, 12pt]{extarticle}
\usepackage[left=0.5in, right=0.5in, top=0.4in, bottom=0.55in]{geometry}
\usepackage{graphicx} % Required for inserting images
\usepackage[normalem]{ulem}
\usepackage{tabto}
\usepackage{minted}
\usepackage{float}
\renewcommand\ULthickness{1.0pt}   %%---> For changing thickness of underline
\setlength\ULdepth{1.3ex}%\maxdimen ---> For changing depth of underline
\usepackage{tikz}
\usetikzlibrary{graphs}

\usepackage{array}
\usepackage{caption}
\usepackage{amsmath} 
\usepackage{fancyhdr}
\usepackage{lastpage}
\usepackage{enumitem}

\author{Juan R Ortiz; Timothy Jung}

% Page style with slide numbers
\pagestyle{fancy}
\fancyhf{}
\renewcommand{\headrulewidth}{0pt}
\fancyfoot[R]{\Large\thepage\ / \pageref{LastPage}}

% Define header command for consistent formatting
\newcommand{\slideheader}[2]{%
\noindent

\includegraphics[height=0.5in]{../../images/cpp_logo.png}%
\hspace{0.3in}%
{\Large \uline{#1 \hfill #2}}%
\par\vspace{0.4in}
}
% Standard content spacing - 0.6in base indent for all content
\newcommand{\contentstart}{%
\leftskip 0.6in \rightskip 0.6in
}
\newcommand{\contentend}{%
\leftskip 0in \rightskip 0in
}

% Standard font size for body content
\newcommand{\bodytext}{\LARGE}

\begin{document}
\begin{huge} 
%%%%%%%%%%%%%%%%%%%%%%%%%%%%%%%%%%%%%%%%
% SLIDE 1 - Title Slide
\slideheader{ECE5110 Numerical Modeling}{Unit 04}

\vspace{1.2in}
\begin{center}
{\Huge \textbf{Juan Ortiz and Timothy Jung}}

\vspace{0.4in}
{\LARGE Solving a 2-Link Robotic Arm Using Jacobian-Based Methods in Python}
\end{center}
%%%%%%%%%%%%%%%%%%%%%%%%%%%%%%%%%%%%%%%%
\newpage %Page 2
\slideheader{Unit 04: Nonlinear Systems}{ECE5110 - Numerical Modeling}

\leftskip 0.7in \rightskip 0.7in 
\vspace{0.4in}
\noindent Robotic arms are widely used in various industries for tasks such as manufacturing, assembly and material handling. 
\vspace{0.4in}

\noindent Our goal is to solve the inverse kinematics of a 2-link robotic arm using Jacobian-based methods in Python.

\begin{center}
\includegraphics[height=3.75in]{../../images/robot_arm.png}%
\end{center}

\vspace{0.2in}

\leftskip 0in \rightskip 0in 

%%%%%%%%%%%%%%%%%%%%%%%%%%%%%%%%%%%%%%%%
\newpage %Page 3
\slideheader{Unit 04: Nonlinear Systems}{ECE5110 - Numerical Modeling}

\leftskip 0.7in \rightskip 0.7in 
\vspace{0.4in}
\noindent The kinematics of a 2-link robotic arm can be expressed as:

\[\hspace{0.4in}
x = L_1\cos(\theta_1) + L_2\cos(\theta_1 + \theta_2) - target_x
\]
\[\hspace{0.4in}
y = L_1\sin(\theta_1) + L_2\sin(\theta_1 + \theta_2) - target_y
\]
\vspace{0.08in}

Where $x$ and $y$ represent the position of the end effector offset from the target position $(target_x, target_y)$, $L_1$ and $L_2$ are the lengths of the robotic arm links, and $\theta_1$ and $\theta_2$ are the joint angles.
\vspace{0.1in}

\noindent The scalar function is defined as:
\[\hspace{0.4in}x_{k+1} = x_k - \frac{f(x_k)}{f'(x_k)} = x_k - \Delta x\]

\vspace{0.4in}


\vspace{0.2in}

\leftskip 0in \rightskip 0in 

%%%%%%%%%%%%%%%%%%%%%%%%%%%%%%%%%%%%%%%%
\newpage %Page 4
\slideheader{Unit 04: Nonlinear Systems}{ECE5110 - Numerical Modeling}

\leftskip 0.7in \rightskip 0.7in 
\vspace{0.2in}
\noindent $\Delta x$ is the change we need to apply to the current guess $x_k$ to get closer to the solution.
\[\hspace{0.4in}J(x_k) \Delta x = -f(x_k)\]
\vspace{-0.2in}
\noindent where:
\[\hspace{0.4in}J(x_k) = 
\begin{bmatrix}
\frac{\partial f_1}{\partial \theta_1} & \frac{\partial f_1}{\partial \theta_2} \\
\frac{\partial f_2}{\partial \theta_1} & \frac{\partial f_2}{\partial \theta_2}
\end{bmatrix} \]

{\huge
\begin{align*}
\frac{\partial f_1}{\partial \theta_1} &= -L_1\sin(\theta_1) - L_2\sin(\theta_1 + \theta_2) \\
\frac{\partial f_2}{\partial \theta_1} &= L_1\cos(\theta_1) + L_2\cos(\theta_1 + \theta_2) \\
\frac{\partial f_2}{\partial \theta_2} &= L_2\cos(\theta_1 + \theta_2) \\
\frac{\partial f_1}{\partial \theta_2} &= -L_2\sin(\theta_1 + \theta_2)
\end{align*}
}

\vspace{0.2in}

\leftskip 0in \rightskip 0in 

%%%%%%%%%%%%%%%%%%%%%%%%%%%%%%%%%%%%%%%%
\newpage %Page 5
\slideheader{Unit 04: Nonlinear Systems}{ECE5110 - Numerical Modeling}

\leftskip 0.7in \rightskip 0.7in 
\vspace{0.2in}
We will use a Jacobian-based Python solver that iteratively updates joint angles.
\begin{center}
\begin{minipage}{0.75\linewidth} % adjust width as you like
\begin{minted}[fontsize=\Large, frame=single]{python}
# Build Jacobian via forward diff.
J = np.zeros((n, n), dtype=float)
for j in range(n):
    x_plus = x.copy()
    x_plus[j] += h
    F_plus = np.asarray(F(x_plus), 
              dtype=float).reshape(-1)
    if F_plus.size != n:
        raise ValueError(
            f"F(x) must return a vector 
              of length {n}, got 
              {F_plus.size}"
        )
    J[:, j] = (F_plus - Fx) / h
\end{minted}
\end{minipage}
\end{center}


\vspace{0.2in}

\leftskip 0in \rightskip 0in 

%%%%%%%%%%%%%%%%%%%%%%%%%%%%%%%%%%%%%%%%
\newpage %Page 6
\slideheader{Unit 04: Nonlinear Systems}{ECE5110 - Numerical Modeling}

\leftskip 0.7in \rightskip 0.7in 
\vspace{0.2in}
The solver approximates the Jacobian using forward finite differences.

\noindent Where:

\[\hspace{0.4in} J^{(k)}_{:.j} \approx \frac{f(x_k + h e_j) - f(x_k)}{h}\]
\vspace{0.4in}

\noindent Leaving us to solve the linear system:

\[\hspace{0.4in}
J(x^{(k)}) \delta^{(k)} = -F(x^{(k)})
\]

\begin{center}
\begin{minipage}{0.75\linewidth} % adjust width as you like
\begin{minted}[fontsize=\huge, frame=single]{python}
delta, err = self.solve_lsoe(J, -Fx)
\end{minted}
\end{minipage}
\end{center}


\vspace{0.2in}

\leftskip 0in \rightskip 0in 

%%%%%%%%%%%%%%%%%%%%%%%%%%%%%%%%%%%%%%%%
\newpage %Page 7
\slideheader{Unit 04: Nonlinear Systems}{ECE5110 - Numerical Modeling}

\leftskip 0.7in \rightskip 0.7in 
\vspace{0.2in}
\noindent Now we can update each step:

\[\hspace{0.4in}
x^{(k+1)} = x^{(k)} + \delta^{(k)}
\]

\begin{center}
\begin{minipage}{0.75\linewidth} % adjust width as you like
\begin{minted}[fontsize=\Huge, frame=single]{python}
x = x + delta
\end{minted}
\end{minipage}
\end{center}



\leftskip 0in \rightskip 0in 

%%%%%%%%%%%%%%%%%%%%%%%%%%%%%%%%%%%%%%%%
\newpage %Page 8
\slideheader{Unit 04: Nonlinear Systems}{ECE5110 - Numerical Modeling}

\leftskip 0.7in \rightskip 0.7in 
\vspace{0.2in}
\noindent Results:

\vspace{-0.2in}
\begin{center}
\includegraphics[height=5.5in]{../../images/unit04_nlsoe_convergence.png}%
\end{center}

\leftskip 0in \rightskip 0in 

%%%%%%%%%%%%%%%%%%%%%%%%%%%%%%%%%%%%%%%%
\newpage %Page 9
\slideheader{Unit 04: Nonlinear Systems}{ECE5110 - Numerical Modeling}

\leftskip 0.7in \rightskip 0.7in 
\vspace{0.2in}
\noindent In this graph, we can see the evolution of the joint angles $\theta_1$ and $\theta_2$ over each iteration of the solver.

\begin{center}
\includegraphics[height=5.0in]{../../images/unit04_nlsoe_state_trajectory.png}%
\end{center}

\leftskip 0in \rightskip 0in 

%%%%%%%%%%%%%%%%%%%%%%%%%%%%%%%%%%%%%%%%

\newpage %End Page
\slideheader{Unit 04: Nonlinear Systems}{ECE5110 - Numerical Modeling}

\leftskip 0.7in \rightskip 0.7in 
\vspace{2.5in}
\begin{center}
{\Huge \textbf{QUESTIONS?}}
\end{center}
\leftskip 0in \rightskip 0in 


%%%%%%%%%%%%%%%%%%%%%%%%%%%%%%%%%%%%%%%%


\end{huge}
\end{document}